\begin{annexitem}{annex:p5.checker-description}{thm:p05-selected-direct-pa-graph}{The specified checker and its internal library}
We first identify the finite descriptions and their numerical values.  We
then prove the particular equations needed inside \(\PA\).  A function
application in an internal argument abbreviates its specified unique graph
output.  Existentially introduce those outputs before the argument and
eliminate them afterwards; \cref{thm:p05-pa-pr-graphs} licenses this for
each fixed description.  Induction below is \(\PA\) induction on the
displayed arithmetic formulas with all other variables retained as parameters.
It does not quantify over descriptions.

\paragraph{Finite expansion and total defaults.}
Predecessor, zero test, multiplication, and truncated subtraction in
\cref{def:p05-checker-description-language} are primitive recursions on their
last arguments.  Sign, equality, order, and conditional selection are their
compositions.  The fold convention inserts one further primitive-recursion
node followed by substitution of its bound.  In particular, the upper bound
of a sum or a search is not an unbounded minimization.  The two search fields
are an ordinary cons-list state; the update is a finite conditional expression
in their projections.  Maps, filters, and the two table schemes expand through
the same operations.

The syntax tables have no circular calls.  Their row expressions consult only
their supplied accumulator; the first table calls the earlier numerical/list
library, and the second may additionally call the already completed first
table.  The substitution and alpha state updates call these two tables and
the list library, never their own completed traversal functions.  Each update
has finitely many cases with fixed command arities.  Substitution is one fold
with bound \(20(n(c)+1)^2\); alpha comparison uses another with bound
\(4(n(c)+1)\).  Prefix stripping and parameter wrapping are subsequent
folds.  Finally the row table calls these earlier functions and consults its
supplied table accumulator.  Its final products and the outer-record case
complete the description of \(\chi_{19}\).  This is a finite acyclic
ordering of description declarations.  Expanding it by the stated projection,
composition, and recursion conventions determines \(d_{19}\).

Every numerical input has a value.  A raw list is always decoded, including
zero.  Wrong syntax gives zero metadata or a failed machine result; an invalid
machine state is absorbing; missing stack values and wrong command arities
fail.  A row failing any required test is exactly \([0,0]\) in the row table.
Bad outer proof tags or arities give zero.  All traversals still use their
fixed bounds on such inputs.  There is consequently no reliance on an
unspecified terminating search even away from well-formed codes.

\paragraph{Identification with the numerical syntax operations.}
The governing pairing is injective and its two inverse components are below
their positive cons code.  Hence in the displayed double sums for \(h,t\)
there is exactly one matching pair when \(c>0\), and none when \(c=0\).
The sums therefore give exactly the earlier head and tail.  The tail reaches
zero by step \(c\); the sum defining length counts precisely its nonzero
tails.  Thus entries, reversal, concatenation, maps, and filters have their
usual specified raw-list values.

Induction over the expression code identifies the first syntax table with the
exact term/formula grammar, node count, and maximum occurring variable index.
Every proper expression child is below its parent.  Induction over the same
code identifies the second table with all and free variable occurrence: a
binder adds its own name to the all-variable field and removes it from the
free-variable field.  The free-index scan therefore gives precisely the
increasing list in \cref{def:p05-canonical-induction-code}.

For the substitution machine, maintain the following execution invariant.
A pending command requests the indicated substitution or free renaming of its
input expression.  Assembly commands retain the exact constructor and binder
that must be wrapped around the returned child values.  The value stack lists
completed child results in reverse return order.  At a binary node, the left
request returns first and the right request second; popping right then left
therefore assembles the original child order.  At a stopped binder the whole
subtree is returned without entering it.  At a conflicting binder the least
fresh search is precisely the search in the deterministic substitution
definition.  Its bound includes one index greater than every forbidden index,
so a permitted index exists.  The renaming command changes exactly the free
occurrences of the old binder name and stops beneath another binder of that
name.  Its new name was absent from the whole matrix, so this pass needs no
additional capture correction.  The continuation then submits that renamed
matrix to the original substitution request and wraps the new binder.
These observations verify each transition against the three quantifier
clauses and the nonbinding clauses of the actual substitution definition.

There are at most \(n(c)\) main node visits: inserted replacement terms are
returned as codes and are never traversed as new main input.  Preliminary
renaming preserves shape.  There are at most \(n(c)\) such renamings, each
visiting at most \(n(c)\) nodes.  A node visit contributes its one call
command and at most one assembly command; a conflicting quantifier contributes
one additional continuation.  Thus at most
\(2n(c)+2n(c)^2+n(c)\) command updates occur on valid input, which is below
the stated \(20(n(c)+1)^2\) bound.  At completion the invariant gives one
value and no pending command.  The absorbing completion state retains this
value for all unused updates.  Structural induction with the transition
invariant consequently identifies the extracted code with the book's exact
substitution, including renaming when the replaced variable is absent freely.

For alpha comparison the environment attached to a traversal command is
exactly the list of its enclosing binder names, nearest first.  The first
matching position therefore codes a bound occurrence by its nearest binder;
failure to find a match codes the original free name.  Quantifier assembly
retains only the quantifier tag and transformed matrix.  All other positions
are unchanged.  There are at most two updates per original node, below
\(4(n(c)+1)\), so the output is exactly the comparison record in
\cref{sec:arithmetic-coding-04}.  Equality of these records is the proved
syntactic alpha test.  This establishes the numerical meanings of both fixed
machines; it asserts no internal truth predicate.

\paragraph{The axiom and row descriptions.}
The \(U\) fold wraps the reversed parameter list.  The substituted base and
step and the constructor expression for \(J\) therefore coincide literally
with the direct induction construction.  Maximal prefix stripping stops at
the first nonuniversal constructor and then retains its state.  On a valid
induction axiom that constructor is an implication; its right child recovers
the induction variable and matrix.  Reconstructing and comparing the complete
code is exactly the final numerical axiom test.  The seven literal constant
tests are unchanged.  The numerical recognition theorem
\cref{thm:p05-direct-pa-axioms} thus identifies \(\mathsf{Ax}_{19}\) with
the direct axiom characteristic function.

Induction through the row list identifies each successful pair
\([1,O_j]\) with the inherited rule check and its list of surviving open
occurrences.  Assumption and identity introduction have the stated empty
premise cases.  Every nondischarge case concatenates its premise occurrence
lists in the inherited order.  Every discharge filter removes precisely the
matching label and literal formula in its specified branch and requires a
strict decrease in length.  Universal introduction tests free occurrences;
existential elimination tests all occurrences after its minor discharge,
including the major open list.  In that latter case the minor remainder is
placed before the major list, exactly as specified in the numerical checker.
All supplied formula/term parameters and literal substitution or alpha checks
are the already identified operations.  A failed local check makes the global
product zero regardless of later rows.  The parent-count product requires one
parent for every row except the last, and zero for the last.  Combined with
earlier references, this is precisely the inherited tree test.

The final root, outer-tag, zero-field, sentence, and surviving-axiom products
are consequently the exact direct proof test.  Thus
\(\chi_{19}(p,x)=1\) iff \(\mathsf{StdProof}_{\PA}(p,x)\).
The specified-graph theorem makes \(\Gamma_{d_{19}}\) a
\(\Sigma^0_1\) formula with exactly these standard values.  Existential
quantification over \(p\), followed by
\cref{thm:p05-standard-pa-numerical-proof}, gives the asserted standard
proof meaning.  This use of standard values has identified the selected
description; the internal claims next use its equations and induction.

\paragraph{Internal head, tail, length, and entries.}
Work now inside \(\PA\).  Its proved pairing and cons equations supply,
for \(c>0\), unique \(a,b<c\) with \(C(a,b)=c\).  Bounded-sum
induction proves that a sum whose only nonzero summand is at index \(a\)
has that summand as its value: split the recursion into indices below the
index, the index itself, and later indices.  Applying this twice to the head
and tail sums gives
\[
 h(0)=t(0)=0,\quad h(C(a,b))=a,\quad t(C(a,b))=b,
 \quad c>0\to c=C(h(c),t(c)).
\]
All uses of equality and sign in these sums are licensed by their zero and
successor equations.  Induction on their recursive arguments proves
\(Z(x)=1\leftrightarrow x=0\),
\(E(x,y)=1\leftrightarrow x=y\), and
\(\mathsf{lt}(x,y)=1\leftrightarrow x<y\); the subtraction step uses
the predecessor equations and the already proved arithmetic order laws.

Induction on \(j\) gives
\[
 T(C(a,b),Sj)=T(b,j),\qquad T(0,j)=0,
 \qquad T(c,j+k)=T(T(c,j),k).
\]
Strong induction on \(c\), available from full induction by retaining all
earlier cases, gives \(T(c,c)=0\).  At positive \(c\), its tail
\(b<c\) satisfies \(T(b,b)=0\) by the induction hypothesis.  Since
\(b\leq D(c)\), the iteration equation and absorption at zero give
\(T(b,D(c))=0\), hence \(T(c,c)=0\).  It follows from the same
iteration equation that \(T(c,j)=0\) whenever \(j\geq c\).

The defining bounded sum for length now yields
\[
 \ell(0)=0,\qquad\ell(C(a,b))=1+\ell(b),\qquad\ell(c)\leq c.
\]
Indeed the zeroth summand at a positive cons is one; reindexing the remaining
summands using the first iteration equation gives the tail's length sum,
possibly followed by zero terms.  Extending a bounded sum by zero terms is
proved by induction on the added interval.  The entry equations are
\[
 e(C(a,b),0)=a,\qquad e(C(a,b),Sj)=e(b,j),\qquad
 j\geq\ell(c)\to e(c,j)=0.
\]
For the last assertion use strong induction on \(c\), the length equation,
and the entry equation at a successor index; the zero-code case is immediate.
These are internal formulas for all numerical codes.

List extensionality also holds internally: equal lengths and equal entries
below that length imply equal codes.  Induct on the common length.  At zero,
the positive-cons length equation forces both codes to be zero.  At a
successor, both codes are positive, their heads agree by the zeroth entry,
and their tails have equal preceding length and entries.  The induction
hypothesis equates the tails, and cons uniqueness equates the lists.

\paragraph{Internal reversal, concatenation, and bounded collection operations.}
The useful invariant for the recursion \(V(c,z,k)\), for
\(k\leq\ell(c)\), is
\[
\begin{aligned}
 \ell(V(c,z,k))&=k+\ell(z),\\
 j<k&\to e(V(c,z,k),j)=e(c,L(D(k),j)),\\
 j<\ell(z)&\to e(V(c,z,k),k+j)=e(z,j).
\end{aligned}
\]
At zero these are its initial-value equation and the entry equations.  At
\(Sk\), the update prepends \(e(c,k)\).  Position zero has that new
value; position \(Sj\) has the preceding value at \(j\).  For \(j<k\),
the arithmetic identity \(L(D(Sk),Sj)=L(D(k),j)\) identifies the reversed
indices; the final segment shifts from \(k+j\) to \(Sk+j\).  The cons
length equation proves the length step.  This verifies the complete invariant
by \(\PA\) induction, or equivalently by
\cref{cor:p05-pa-recursion-invariant} with these base and step proofs.

Apply it first to \(\mathsf{rev}(c)\), then to
\(V(\mathsf{rev}(c),d,\ell(c))\).  The double reversal of an index
\(j<\ell(c)\) is \(j\), by subtraction of ordered pairs and addition
cancellation.  We obtain
\[
\begin{aligned}
 \ell(\mathsf{cat}(c,d))&=\ell(c)+\ell(d),\\
 j<\ell(c)&\to e(\mathsf{cat}(c,d),j)=e(c,j),\\
 j<\ell(d)&\to e(\mathsf{cat}(c,d),\ell(c)+j)=e(d,j).
\end{aligned}
\]
The zero cases and associativity of concatenation follow by these equations
and list extensionality, dividing the possible entry indices into their
successive intervals.  This also proves that appending an entry preserves
all earlier positions and places the new entry at the old length.

For a fixed map expression, induction on its fold counter gives length equal
to that counter and the mapped value at each earlier input position.  Its
step is the just-proved append equation.  At the final counter the map
therefore preserves length and has exactly the declared entries.  For a
filter, the analogous induction says that its result consists, in their
original order, of the earlier entries whose bit was one.  At a step either
that next entry is appended or the entire state is retained.  In particular,
every result entry is an input entry and a strict length decrease is
equivalent to removal of at least one input position.  These statements may
be expressed without the word ``consists'' by the increasing list of retained
input positions, constructed by the same filter-counter fold.  Its entries
give the position witnesses; the two update cases prove the assertions.

Two consequences will be used below.  Mapping commutes with concatenation:
use the map-entry equations on the two index intervals and extensionality.
Further, if \(s(j)=a+j\) and two filter tests have equal values on
corresponding input entries, then
\[
 \mathsf{filter}_{b'}(\mathsf{map}_s(O))
       =\mathsf{map}_s(\mathsf{filter}_b(O)).
\]
Induct over the filter counter.  The equal bits select the same branch;
retention preserves the equality, and appending corresponding entries
preserves it by the map/append equation.  Length and nonvacuity are preserved.
Finally, a bounded product of zero-one factors is one exactly when all of
its factors are one: the empty product is one and the successor step uses
the two possible values of its last factor.  Its concatenation law follows
by the same induction and the concatenation-entry equations.

\paragraph{Internal table and syntax equations.}
For each fixed table update, induction on its counter proves that its table
has that length.  Appending a row leaves all earlier entries unchanged.
Induction on a larger counter therefore proves that the table at an earlier
counter is its corresponding prefix; the defining recursion identifies the
new row with the update applied to exactly that prefix.  Consequently
retrieving child-code metadata from a later table gives the same values as
retrieving them from their own completed prefix tables.  This is an internal
prefix statement about the particular recursion, proved using the append
equations, with no quantification over algorithms.

Constructor injectivity and the internal head/tail equations give all fixed
arity decodings.  A proper expression child is below its enclosing cons
code: each is below its containing argument-list cons, whose tail and head
are themselves below the next cons, ending at the constructor.  Therefore
the first syntax-table update really consults earlier rows at valid nodes,
also inside \(\PA\).  Its fixed Boolean cases and the prefix equations give,
in particular,
\[
\begin{aligned}
 \mathsf{Fm}(N_{10}(u,v))&=\mathsf{Fm}(u)\mathsf{Fm}(v),\\
 \mathsf{FV}(N_{10}(u,v),i)&=
        \mathsf{FV}(u,i)\lor\mathsf{FV}(v,i)
            \quad\text{if both children are formulas}.
\end{aligned}
\]
The second equation follows by the identical prefix argument for the
variable table.  Strong induction on an expression code additionally gives
\[
 i\geq c\quad\longrightarrow\quad
             \mathsf{AV}(c,i)=\mathsf{FV}(c,i)=0.
\]
At an invalid code this is its default.  At a variable its index is below
the variable code.  At a compound every child code is smaller, so the
induction hypotheses kill its fields; at a quantifier the binder index is
also below the code.  Its Boolean row equations then give zero.  This proves
the free-index bound internally rather than importing an external syntax fact.

The filter invariant for \(F(c)\) now identifies its empty-list test with
\(\forall i<c\,\mathsf{FV}(c,i)=0\) when \(c\) is a formula.
Suppose \(\mathsf{Sent}(N_{10}(u,v))=1\).  Both children are formulas
by the first displayed equation.  For every \(i<u\), and for every
\(i<v\), that index is below the implication code.  The zero free-variable
value at that index, together with the Boolean disjunction equation, forces
the corresponding child's value to be zero.  The two free-index folds are
therefore empty.  This proves \(\mathsf{Sent}(u)=\mathsf{Sent}(v)=1\)
inside \(\PA\), including children with no free variables initially.

All list, table, and syntax assertions used later have thus received their
base and step proofs for the selected graphs.  Internal semantic properties
of the substitution or alpha machines are not needed for the joining proof:
there their complete numerical arguments are literally unchanged, so unique
outputs of the same fixed descriptions suffice.  The concrete identity-proof
example in the body has one successful identity-introduction row, zero parent
count, a sentence root, and an empty final axiom product.  These particular
values, by finite use of the graph equations on their numerals, prove its
stated \(\PA\) instance.
\end{annexitem}

\begin{annexitem}{annex:p5.checker-internal-mp}{thm:p05-direct-pa-internal-mp}{The internal proof-joining invariant}
Work inside \(\PA\) with the unique graph-output conventions and library
equations proved in \cref{annex:p5.checker-description}.  Every invariant
below is an arithmetic formula: list membership means a bounded entry
quantifier, and every function application expands through its fixed graph.
We prove the assertions for all numerical input codes satisfying the checker
premises, including codes that are nonstandard in a model of \(\PA\).

\paragraph{The offset equations.}
The map invariant for \(\mathsf{sh}\) gives
\[
 \ell(\mathsf{sh}(a,J))=\ell(J),\qquad
 j<\ell(J)\to e(\mathsf{sh}(a,J),j)=a+e(J,j).
\]
Addition cancellation makes this map injective on values, and preservation
of strict order by addition gives
\(i<j\leftrightarrow a+i<a+j\).  The definition of \(\mathsf{sr}\)
and fixed cons decoding show that a valid row's rule tag, conclusion formula,
parameter list, and marker are unchanged; only its premise list has the
displayed offset.  The map invariant for \(\mathsf{SR}\) gives its original
row count and that shifted row at each original position.

For row lists \(Q,P\) of lengths \(a,b\), set temporarily
\[
 R=\mathsf{cat}(Q,\mathsf{SR}(a,P)).
\]
Concatenation entries give the complete lookup equations
\[
\begin{aligned}
 i<a&\to e(R,i)=e(Q,i),\\
 j<b&\to e(R,a+j)=\mathsf{sr}(a,e(P,j)),\\
 j<b&\to f_R(a+j)=f_P(j),\quad
                       \lambda_R(a+j)=\lambda_P(j).
\end{aligned}
\]
The last line assumes the original rows have length five, as they do under
the old checker premises.  Their parameter lists and formula/term witnesses
have the same unchanged-field equation.  These equations also hold after
one final row is appended to \(R\), since the append preserves its prefix.

\paragraph{The open-occurrence invariant.}
First consider a row list all of whose row bits are one.  Induction on its
row counter proves that every entry of \(O_j\) is an assumption-row index
at most \(j\).  At an assumption the list is \([j]\).  Every premise
index is smaller; its list satisfies the induction hypothesis.  Nondischarge
rules concatenate those lists, and discharge rules remove entries from the
specified lists.  Identity introduction supplies the empty list.  These
cases prove the assertion, using the list/filter equations rather than a
metatheoretic interpretation of the rows as a tree.

Suppose all row bits of \(Q\) and \(P\) are one.  Induct over the first
\(a\) row updates for \(R\).  Its row, every earlier reference, every
formula/marker lookup at an open index, and every already computed premise
open list agree with those for \(Q\).  The fixed row expression consequently
has the same bit and same open list.  This proves the first block invariant.

For the second block prove by induction on \(j<b\) that its bit is the old
bit one and its open list is
\[
             O^{R}_{a+j}=\mathsf{sh}(a,O^{P}_j).
\]
Every premise reference \(k<j\) becomes \(a+k<a+j\).  The preceding
offset lookup equations give identical premise formulas and parameter lists;
the induction hypothesis supplies the shifted premise open lists.  To check
the rule update, the following is an exhaustive partition of its cases.
\begin{enumerate}
\item At an assumption, the new list is \([a+j]\), which is the shift
 of \([j]\); all its remaining tests are unchanged.  At identity
 introduction, the term parameter and identity conclusion are unchanged and
 both open lists are empty.
\item The nondischarge propositional rules, universal elimination,
 existential introduction, alpha conversion, and identity elimination use
 only the unchanged formula and parameter data for their additional tests.
 Every substitution or alpha call therefore has exactly the same arguments.
 Its specified graph is single-valued, so its computed output is equal in
 the two computations.  Concatenation commutes with shift, as proved for
 maps, giving the new open-list equation.
\item For implication introduction, negation introduction, and RAA, the
 discharged formula and label are unchanged.  At every old open index \(k\),
 the new index \(a+k\) has exactly its old formula and marker.  The two
 removal bits are equal.  The filter/shift equation therefore gives the
 shifted filtered list, and preserved lengths give exactly the same
 nonvacuity test.  Disjunction elimination applies this argument separately
 to its second and third premise branches; its major list is preserved by
 concatenation.  Labels that coincide in the two copied blocks cannot cause
 any extra removal, since an old discharge consults only its named premise
 open lists.
\item At universal introduction, the generalized variable and all open-leaf
 formulas are unchanged.  Thus the \(\mathsf{FV}\) calls at corresponding
 positions have identical arguments and outputs.  Equal lengths and the
 bounded-product induction give the same freshness test and shifted list.
\item At existential elimination, its matrix, bound-variable index, witness
 variable, and discharge label are unchanged.  Hence its computed minor
 assumption formula is unchanged.  The filter argument gives exactly the
 shifted minor remainder and preserves nonvacuity.  Concatenating that
 remainder with the major list gives the shifted remaining open list in the
 required order.  The conclusion and major formula are unchanged, as is
 every leaf formula in this list.  Each of its all-variable calls therefore
 has the same arguments; its three factors and its bounded product are
 equal.  This proves the actual all-occurrence freshness test, including
 bound occurrences, for the copied row.
\end{enumerate}
The common arity tests and zero-marker conditions are unchanged in every
case, and the addition-order equation already established the earlier-reference
test.  These cases exhaust tags zero through nineteen and prove both parts of
the second block invariant.  Appending a last row changes no lookup used in
the two block computations, so the same proof applies to their prefix in
the completed joined record.

\paragraph{The parent-count invariant.}
Assume additionally \(\mathsf{Tree}(Q)=\mathsf{Tree}(P)=1\).
Then \(a,b>0\), and the zero-one product lemma supplies all the row-bit
hypotheses used above.  Let
\[
 R'=\mathsf{cat}(R,[[8,v,[D(a),D(a+b)],0,0]]).
\]
Every reference from an old \(Q\) row is below \(a\); every reference
from a shifted \(P\) row has the form \(a+k\) with \(k<b\).
Split the defining outer sum for \(\mathsf{pc}\) into these two intervals
and the last row, using bounded-sum induction and the concatenation-entry
equations.  Within the second interval use
\(E(a+k,a+j)=E(k,j)\), from cancellation.  No reference in one old block
can contribute to a position in the other.  The two last-row references
are distinct, because \(D(a)<a\leq D(a+b)\) when \(a,b>0\).
We obtain
\[
\begin{aligned}
 i<a&\to\mathsf{pc}(R',i)=\mathsf{pc}(Q,i)+E(i,D(a)),\\
 j<b&\to\mathsf{pc}(R',a+j)=\mathsf{pc}(P,j)+E(j,D(b)),\\
 \mathsf{pc}(R',a+b)&=0.
\end{aligned}
\]
For the last equality every old reference is below its old row, and the two
new references are below \(a+b\).  The old parent-count tests state that
each old nonroot has count one and each old root count zero.  Substitution
into these equations gives count one at every position below \(a+b\),
and zero at that final position.  This proves the exact new parent-count
product.  It does not assume an external finite-tree interpretation.

\paragraph{The last row and direct axiom tests.}
Assume now the arithmetic checker premises
\[
 \chi_{19}(p,N_{10}(u,v))=1,
                \qquad \chi_{19}(q,u)=1.
\]
Unfold their characteristic descriptions.  The zero-one product property
gives the exact tag-19 decodings, old tree tests, root equalities, sentence
tests, and all old direct-axiom factors.  Put \(Q=R_q,P=R_p\) in the
preceding arguments.  The two old roots occupy positions \(D(a)\) and
\(D(a+b)\); their formulas are \(u\) and \(N_{10}(u,v)\).
The appended row has tag 8, those references in that order, zero parameters,
zero marker, and conclusion \(v\).  Both references are earlier than its
position.  The implication-subformula sentence lemma gives
\(\mathsf{Sent}(v)=1\), and hence \(\mathsf{Fm}(v)=1\).
Its common and additional row tests are therefore one.  It discharges nothing,
and the open-occurrence invariant gives exactly
\[
 \mathsf{open}(R')=
       \mathsf{cat}(\mathsf{open}(Q),
                         \mathsf{sh}(a,\mathsf{open}(P))).
\]

At a leaf from the first segment its formula is exactly the old \(Q\)
formula.  At a leaf \(a+j\) from the second segment its formula is exactly
\(f_P(j)\).  Consequently each new
\(\mathsf{Ax}_{19}\) call has exactly the same numerical argument as its
old successful call.  Unique outputs of that specified graph give equal
values.  The product/concatenation and map-entry equations prove that the
entire new axiom product is the product of the two old products, hence one.
The empty-list case uses the initial product one and is included.  No
internal certificate search, axiom completeness assertion, or replacement
of an axiom by a theorem is used here.

Together with the parent-count result and the last-row check, this proves
that the explicitly assembled \(N_{19}(N_{15}(R'),0)\) passes
\(\chi_{19}(\cdot,v)\).  The success guard of \(\mathsf{mp}_{19}\)
holds by the two checker premises and constructor injectivity, so its
composition and case equations identify this assembled value with its
specified graph output.  We have proved inside \(\PA\), with universal
closure in \(p,q,u,v\),
\[
\begin{split}
 &\chi_{19}(p,N_{10}(u,v))=1\land\chi_{19}(q,u)=1\\
 &\quad\longrightarrow\exists r\,
       (\Gamma_{d_{\rm mp}}(p,q,r)\land\chi_{19}(r,v)=1).
\end{split}
\]
Here every characteristic-function occurrence means exactly its
\(\Gamma_{d_{19}}\) output.  Substituting the definition of
\(\mathsf{Prf}^{19}_{\PA}\) is therefore syntactic graph expansion,
not an inference from standard-model agreement.  This is the first assertion
of \cref{thm:p05-direct-pa-internal-mp}.

\paragraph{The derivability conditions actually obtained.}
For fixed \(\varphi,\psi\), their implication code is the numerical
constructor \(N_{10}(\ulcorner\varphi\urcorner,
\ulcorner\psi\urcorner)\).  Finite evaluation of its specified description
on these numerals gives the corresponding equality in \(\PA\).
Instantiate the internal joining theorem at those two numeral codes.
From the two existential proof assertions, choose \(p,q\), apply that
theorem to obtain \(r\), and existentially introduce the concluding proof
assertion.  Implication introduction yields the stated D2 formula.  Bound
variables are chosen fresh in these finite first-order steps.

For D1, if \(\PA\derives\sigma\), the numerical proof-meaning theorem
gives a standard direct proof code \(p\).  Thus the fixed description has
the actual standard value \(\chi_{19}(p,\ulcorner\sigma\urcorner)=1\).
The numeral-computation clause of \cref{thm:p05-pa-pr-graphs}, applied to
this description and tuple, supplies a \(\PA\) derivation of the selected
graph at output \(\bar1\).  Existential introduction supplies
\(\mathsf{Prv}^{19}_{\PA}(\overline{\ulcorner\sigma\urcorner})\).
This is a metatheoretic implication using a finite standard computation.
The D2 argument above was a universal internal invariant, a different claim.

This D1/D2 argument is distinct from converting bounded evidence into an
annotated proof.  The original-sentence descriptions
\(\widehat d_\sigma^+\) and their evidence-to-proof invariant are supplied
in \cref{thm:p05-original-sigma-evidence}; applying that construction to
the literal selected graph gives D3 in \cref{thm:p05-direct-pa-d3}.
Neither the present proof nor those later constructions equate the selected
graph internally with the earlier arbitrary numeralwise representative
in \(\mathsf Q\).
\end{annexitem}
